\section{Introduction}\label{sec:intro}

Interchange interventions test whether a network's internal state carries the variables an explanation assigns to
it. One runs the network on a recipient input, replaces the part of an internal state that is said to carry a
variable with the corresponding part from a donor input, and checks that the output changes as the explanation
predicts \citep{vig2020gender,geiger2021causal,wang2023ioi}. Causal abstraction makes the underlying claim exact:
the explanation must agree with the network under every intervention on its variables, including interventions on
several variables at once \citep{beckers2019abstracting,geiger2025abstraction}. In practice explanations are
mostly scored one variable at a time, for example by single-variable interchange accuracy
\citep{huang2024ravel,mueller2025mib}. Recent benchmarking work argues that interventions on several variables
catch interactions that single checks miss \citep{meloux2026validating}. So when can a joint intervention reject
an explanation that single interventions accept?

We give a necessary condition for a common kind of explanation, an alignment that maps two variables to orthogonal
linear subspaces at one site, as distributed alignment search produces \citep{geiger2024das}. A joint interchange
edit is exactly a single edit applied to the result of the other single edit (\cref{lem:composition}). If the
single edits of either variable reproduce the natural states of the inputs they simulate, every joint query is an
ordinary single query, and a joint test cannot reject what single tests accept (\cref{prop:redundant}). A joint test
can do so only where the single edits of both variables give the correct output without reproducing those natural
states (\cref{cor:irreducible}): their effect is a dormant behavioural change in the sense of
\citet{grant2026divergent}, invisible in a single test and revealed by the second edit. A small explicit network
shows that this can happen. A quantitative version (\cref{prop:margin}) says how large the misses must be: a joint
query differs from each single query it reduces to by part of the other edit's miss, and it can fail only if, in both
orderings, that difference lowers the output's margin by more than the single query's own margin.

We then examine four settings in which every single and joint intervention can be evaluated, because each task has
a small finite domain. Study~1 revisits a frozen comparison of eight test panels on \CmpPrograms{} compiled
transformers with exact validity labels, whose joint tests had shown no advantage over full-domain single tests.
Study~2, frozen before its final networks were trained, repeats the analysis on \TrnNetworks{} small trained
transformers and MLPs, where linear alignments are approximate. Study~3, whose two predictions were recorded before
any of its networks existed, places the two variables at different sites, where the identity holds in one direction
only. Study~4 asks whether the results hold in a language model: GPT-2 Small on indirect-object identification
\citep{wang2023ioi}, with the two variables of the MIB benchmark's model of the task \citep{mueller2025mib}, the
indirect object's position and its identity. Its protocol and predictions were hashed before any model was run, its
settings were chosen on development names and templates, and its confirmation used held-out ones; an amendment,
recorded before any of its computation, repeated the one-site arm in GPT-2 Medium. In all four studies we compute
exactly how likely a test of a given budget is to detect a failing alignment.

\paragraph{Findings.}
\begin{itemize}
  \item \emph{A necessary condition.} For orthogonal alignments at one site, a joint test can reject what single
  tests accept only if both variables' single edits give correct outputs yet miss
  the natural states they simulate, and only where the misses are large relative to the single queries' margins
  (\cref{sec:theory}).
  \item \emph{In compiled transformers, no alignment could fail only jointly.} Study~1's alignments are by
  construction either exact or wrong on single queries. A post hoc census of all \CmpCandidates{} confirms this and
  explains the frozen comparison's null result within sampling error; it cannot test the condition. The failures
  that restricted single tests missed came from alignments fitted on part of the input range (\cref{sec:compiled}).
  \item \emph{In small trained networks the condition holds wherever both variables change, yet joint-only failures
  are rare.} In Study~2 no single edit that changes a variable reproduces a natural state exactly. Of the
  \TrnPerfect{} alignments (out of \TrnCandidates) that pass every single query, \TrnPerfectJointFailWord{} fail a
  joint query, each on fewer than 1\% of joint queries. Post hoc, the signed first-order margin rule, computed from
  single queries alone, flags \TrnMbRuleFlagged{} of the seven and no other alignment (\cref{sec:trained}).
  \item \emph{There, at the prespecified costs, joint tests detect fewer failing alignments than single tests.} With
  three calls per single and four per joint query, joint minus single detection probability is
  \TrnContrastOneTwenty{} \TrnContrastOneTwentyCI{} at 120 calls and \TrnContrastTwelve{} at 12; after a single test
  of 120 calls, 120 more calls detect a further \TrnAddonJoint\% of the full-domain DAS alignments if spent on joint
  queries and \TrnAddonSingle\% if spent on single queries (post hoc; \cref{sec:trained}).
  \item \emph{With the variables at different sites, joint tests did slightly better, for a simple reason.} Both
  predictions of Study~3 held on their estimates, and joint minus single detection at 120 calls is \MsContrast{}
  \MsContrastCI. Post hoc, most of the gain comes from joint queries always exercising the downstream variable, which
  was the weaker one. The margin rule, fixed on Study~2 and applied without change, flags exactly the
  \MsPerfectJointFailWord{} alignments that fail only jointly (\cref{sec:multisite}).
  \item \emph{In GPT-2, the prediction carried over from the small networks failed.} At one site, joint tests detect
  far more failing alignments than single tests drawn from the full-domain panel, by \ContrastS{} in GPT-2 Small and
  \ContrastM{} in GPT-2 Medium at 120 calls. The position subspace found by distributed alignment search also carries
  the names: it fails \UnchRTwoS\% (rank 2) and \UnchRFourS\% (rank 4) of the single position queries whose donor keeps
  the recipient's position, queries that the full-domain panel leaves out; every joint query carries such a donor.
  Against single tests that also draw these queries, joint tests still detect more, by \ContrastAllS{} and
  \ContrastAllM, a comparison added after the development results (\cref{sec:gpt2-results}).
  \item \emph{In GPT-2, alignments that pass every single query are rare, and those that exist fail jointly.} None does
  in GPT-2 Small. In GPT-2 Medium \PerfectWordM{} do, too few to score the prespecified predictions; all
  \PerfectWordM{} fail joint queries, and the margin rule flags all of them.
  \item \emph{Under the counterfactual model, MIB's token head fails almost every token query} (\HeadsRateB\% at head
  8.10), and the final position alone fails about half of all single queries.
  \item \emph{Subspaces fitted one variable at a time make joint-only failures more common} in Study~2's networks
  (\DShare\% against \DRefShare\%), and their overlap predicts which; in GPT-2 the overlap is small, and joint failures
  fall slightly as it grows.
  \item \emph{At small budgets, guiding joint queries by the margin rule does not pay} in Studies~2 and~3 (post hoc;
  \cref{sec:strategies}).
\end{itemize}

The results concern linear alignments at one site (two in Study~3), finite tasks with two variables, small networks
and two GPT-2 models on one task. \Cref{sec:discussion} discusses what they imply for testing alignments.
